% --- Back cover ---
% Large PERSONIX wordmark, no logo image. Cream paper inherited globally.
% Absolute (tikz overlay) positioning so the footer sticks to the bottom.
\clearpage
\thispagestyle{empty}
\begin{tikzpicture}[remember picture, overlay]
  % --- Wordmark ---
  \node[anchor=north, font=\sffamily\bfseries\fontsize{48}{52}\selectfont]
    at ([yshift=-26mm]current page.north) {PERSONIX};
  \node[anchor=north, font=\rmfamily\itshape\large]
    at ([yshift=-44mm]current page.north) {Անզիջում փոփոխություն};
  \node[anchor=north, font=\rmfamily\small,
        text width=\paperwidth-50mm, align=center]
    at ([yshift=-54mm]current page.north)
    {Անգրաքննելի և անկաշառելի ապակենտրոնացված հեղինակության ցանց};
  % --- Body ---
  \node[anchor=north, text width=\paperwidth-50mm, align=justify, font=\small]
    at ([yshift=-72mm]current page.north)
    {%
      Պետությունն աշխատում էր, երբ հաղորդակցությունը դանդաղ էր, որոշումները՝
      տեղական, իսկ հեղինակությունն ապրում էր հենց այն քաղաքում, որը կառավարում էր։
      Այսօրվա ցանցերը շրջում են այս երեք ենթադրություններն էլ --- և դրանց վրա
      կառուցված ինստիտուտներն այլևս չեն համապատասխանում իրենց կառավարած աշխարհին։
      Այս գիրքը նկարագրում է գործիք, որը համապատասխանում է․
      ապակենտրոնացված հեղինակության ցանց, որտեղ ինքնությունը քոնն է ու մնում է քեզ,
      ճշմարտությունը հրապարակայնորեն ստուգելի է, իսկ կաշառքը հեղինակության ինքնասպանություն է։

      \smallskip
      Ոչ մանիֆեստ։ Ոչ կանխատեսում։ Աշխատող նախագիծ այն մասին, թե ինչպես կարելի է
      կառուցել պետության իրավահաջորդը --- կամավոր, մեկ հայտարարություն մյուսի հետևից։
    };
  % --- Footer (anchored to the bottom of the page) ---
  \node[anchor=south, font=\sffamily\footnotesize,
        text width=\paperwidth-50mm, align=center]
    at ([yshift=12mm]current page.south)
    {%
      Pavel Kudrna\quad\textbullet\quad Պրահա, 2026\par
      \href{https://personix.org}{personix.org}\par
      Արտոնագրված է Creative Commons Attribution-ShareAlike 4.0 International
      (CC BY-SA 4.0) լիցենզիայով։
    };
\end{tikzpicture}%
\mbox{}%
\clearpage
